\documentclass[a4paper,11pt]{article}
\usepackage[utf8]{inputenc}
\usepackage[T1]{fontenc}
\usepackage[italian]{babel}
\usepackage{lmodern}
\usepackage{geometry}
\usepackage{longtable}
\usepackage{array}
\usepackage{booktabs}
\usepackage{xcolor}
\usepackage{pdflscape}
\usepackage{ragged2e}
\usepackage{microtype}

\geometry{
  a4paper,
  top=2cm, bottom=2cm,
  left=1.5cm, right=1.5cm
}

\setlength{\LTcapwidth}{\textwidth}

\newcolumntype{P}[1]{>{\RaggedRight\arraybackslash}p{#1}}

\begin{document}

\appendix
\section*{Appendice -- Risposte baseline e risposte del sistema per topic di valutazione}
\addcontentsline{toc}{section}{Appendice -- Risposte baseline e risposte del sistema}

\noindent
La tabella seguente riporta, per ciascuno dei 15 topic del set di valutazione, la risposta baseline (generata manualmente come riferimento) e la risposta prodotta dal sistema RAG parlamentare.
\bigskip


\begin{landscape}
\small
\setlength{\tabcolsep}{6pt}
\renewcommand{\arraystretch}{1.3}

\begin{longtable}{@{}P{3.5cm} P{9.5cm} P{9.5cm}@{}}
\toprule
\textbf{Topic} & \textbf{Risposta Baseline} & \textbf{Risposta del Sistema} \\
\midrule
\endfirsthead

\multicolumn{3}{c}{\small\itshape (continua dalla pagina precedente)} \\[4pt]
\toprule
\textbf{Topic} & \textbf{Risposta Baseline} & \textbf{Risposta del Sistema} \\
\midrule
\endhead

\midrule
\multicolumn{3}{r}{\small\itshape (continua nella pagina successiva)} \\
\endfoot

\bottomrule
\endlastfoot

Aiuti militari all'Ucraina & \#\# Introduzione\\[6pt] Il dibattito parlamentare in esame verte sulle misure urgenti per la proroga dell'autorizzazione alla cessione di mezzi, materiali ed equipaggiamenti militari in favore delle autorità governative dell'Ucraina, affrontate nei disegni di legge S. 389, S. 974 e S. 1335. Le discussioni si estendono dal novembre 2022 al gennaio 2026, comprendendo gli interventi di oltre trenta deputati in svariate assemblee, tra cui le sedute 15, 240 e 594. I documenti delineano il quadro normativo ed esecutivo sull'impegno dell'Italia rispetto al conflitto in corso senza fornire valutazioni o anticipare gli schieramenti.\\[6pt] \#\# Posizione del Governo\\[6pt] Per il Governo, la Presidente del Consiglio \textbf{Meloni} (Fratelli d'Italia, 17/12/2025) evidenzia che «in gioco non ci sono solo la dignità, la libertà e l'indipendenza dell'Ucraina, ma anche la sicurezza dell'Europa nel senso più ampio del termine». La Premier garantisce inoltre l'impegno dell'esecutivo a mantenere un fronte compatto a livello europeo per ostacolare attivamente le risorse finanziarie che alimentano la macchina bellica della Federazione Russa.\\[6pt] \#\# Posizioni della Maggioranza\\[6pt] Per Fratelli d'Italia, \textbf{Caiata} (Fratelli d'Italia, 30/11/2022) rileva che «il nostro impegno, oltre a essere un sostegno militare, deve innanzitutto essere un sostegno umanitario». Il deputato ritiene fondamentale limitare gli effetti della crisi per la popolazione e rinforzare i legami interni all'Alleanza Atlantica per giungere alla fine delle ostilità.\\[6pt] Per la Lega - Salvini Premier, \textbf{Zoffili} (Lega - Salvini Premier, 15/01/2026) dichiara che «la Lega, e lo ha detto anche il Ministro \textbf{Crosetto}, non accetterà mai l'invio di militari italiani in Ucraina». Il gruppo difende la fornitura di soli sistemi difensivi orientati a proteggere la popolazione civile e le infrastrutture sensibili dai raid missilistici.\\[6pt] Per Forza Italia, \textbf{Battilocchio} (Forza Italia, 23/01/2023) ribadisce che «è il momento di non cedere e di non esitare, confermando, sia l'aiuto umanitario, sia quello militare». Il parlamentare insiste sulla necessità di supportare costantemente Kiev per impedire una resa, in attesa che si concretizzi un valido e giusto tavolo negoziale.\\[6pt] Per Noi Moderati, \textbf{Romano} (Noi Moderati, 15/01/2026) nota che «non esiste un negoziato serio se una delle parti è costretta a trattare sotto i bombardamenti». La forza politica giudica l'assistenza armata come presupposto fondamentale e ineludibile per poter avviare la diplomazia in condizioni di reale equilibrio.\\[6pt] \#\# Posizioni dell'Opposizione\\[6pt] Per il Partito Democratico, \textbf{Porta} (Partito Democratico, 07/02/2024) ritiene che «nessuno di noi crede che l'Ucraina non debba esercitare il suo diritto a difendersi e che la comunità internazionale, Europa in primis, non possa aiutare con tutti i mezzi possibili l'eroica resistenza del popolo ucraino». La fazione chiarisce che gli sforzi tattici sul campo devono progredire di pari passo con un'attiva e autorevole iniziativa di pace.\\[6pt] Per il Movimento 5 Stelle, \textbf{Pellegrini} (Movimento 5 Stelle, 22/01/2025) denuncia che «se la fornitura di armi non è accompagnata da una penetrante, convincente, direi quasi ossessiva attività diplomatica e negoziale, l'unico risultato di inviare le armi sempre più potenti è di alimentare la guerra, non certo di circoscriverla». Il partito esorta il Parlamento a bloccare ulteriori pacchetti bellici per tentare di superare l'attuale stallo internazionale.\\[6pt] Per Alleanza Verdi e Sinistra, \textbf{Zanella} (Alleanza Verdi e Sinistra, 23/01/2023) contesta che «l'azione di questo Governo, sulla scia di quello precedente, si sta limitando solo alle forniture militari». Il gruppo condanna la carenza di sforzi tangibili per ottenere il cessate il fuoco e critica fermamente i profitti miliardari accumulati dalle industrie delle armi.\\[6pt] Per Azione, \textbf{Rosato} (Azione, 28/02/2024) afferma che «purtroppo la pace, la libertà e la democrazia si difendono proprio così, anche con le armi». Il raggruppamento sollecita le istituzioni nazionali ed europee ad accelerare la produzione di artiglieria per bilanciare l'economia di guerra russa e proteggere l'intero continente.\\[6pt] Per Italia Viva, \textbf{Del Barba} (Italia Viva, 15/10/2024) propone che «accanto a questa immediata condanna, siano seguiti gli aiuti su più fronti, compreso quello delle armi». Il deputato insiste affinché l'Unione Europea non arretri e invoca l'utilizzo diretto delle ingenti risorse russe congelate a garanzia dei fondi per la difesa ucraina.\\[6pt] Per il Misto, \textbf{Della Vedova} (Misto, 08/02/2024) precisa che «è ancora più importante garantire agli ucraini la possibilità di difendersi dall'aggressione russa, perché \textbf{Putin} scommette esattamente sulla nostra fatica». L'esponente invita a impiegare legalmente gli immensi asset finanziari sottratti alla Federazione Russa in Occidente come contromisura per finanziare l'assistenza. & \\
\midrule
Autonomia Differenziata & \#\# Introduzione\\[6pt] Il dibattito parlamentare sull'attuazione dell'Autonomia Differenziata si concentra principalmente sul disegno di legge S. 615 (A.C. 1665) e sulle relative mozioni, sviluppandosi in un periodo temporale che va dal dicembre 2022 al settembre 2024. Le discussioni hanno coinvolto oltre venti deputati in numerose sedute dell'Assemblea, tra cui spiccano le sedute 285 e 309, specificamente dedicate all'esame dettagliato del provvedimento. L'ampia e costante partecipazione documenta la complessa rilevanza politica della materia e le profonde divergenze tra gli schieramenti.\\[6pt] \#\# Posizione del Governo\\[6pt] Per il Governo, il Ministro per gli Affari Regionali e le Autonomie \textbf{Calderoli} (Lega - Salvini Premier, 18/06/2024) afferma che «l'autonomia differenziata è nella nostra Costituzione dal 2001; è stata approvata dal Parlamento e da un referendum popolare». Il rappresentante dell'esecutivo difende la coerenza della norma costituzionale, sostenendo fermamente che le funzioni decentrate permettono di erogare i servizi in modo molto più efficiente rispetto alla gestione statale centralizzata. \\[6pt] \#\# Posizioni della Maggioranza\\[6pt] Per Fratelli d'Italia, \textbf{Gardini} (Fratelli d'Italia, 29/04/2024) rileva che «l'autonomia differenziata non è un'improvvisata, ma è una previsione costituzionale, approvata più di 20 anni fa, della quale regioni di centrodestra e di centrosinistra hanno chiesto nel tempo l'attuazione». La deputata difende l'impianto della legge rimarcando la presenza di clausole di garanzia e il superamento del criterio della spesa storica per tutelare tutte le aree del Paese.\\[6pt] Per la Lega - Salvini Premier, \textbf{Furgiuele} (Lega - Salvini Premier, 29/04/2024) evidenzia che «l'autonomia differenziata è un'opportunità, una scelta libera e non un obbligo per le regioni». L'esponente del partito assicura che il provvedimento garantirà responsabilità, meritocrazia e sburocratizzazione, tutelando sia i territori richiedenti sia quelli che sceglieranno di non avviare le intese.\\[6pt] Per Forza Italia - Berlusconi Presidente - PPE, \textbf{Pagano} (Forza Italia - Berlusconi Presidente - PPE, 18/06/2024) auspica un solido decentramento amministrativo chiarendo che l'autonomia «non serve a far altro che ad avvicinare le istituzioni ai cittadini, avvicinare, soprattutto, il momento della decisione ai cittadini». Il gruppo supporta la valorizzazione delle peculiarità territoriali, insistendo sulla necessità di gestire il processo con equilibrio per evitare nuove disuguaglianze.\\[6pt] Per Noi Moderati (Noi con l'Italia, Coraggio Italia, UDC e Italia al Centro) - MAIE - Centro Popolare, \textbf{Semenzato} (Noi Moderati (Noi con l'Italia, Coraggio Italia, UDC e Italia al Centro) - MAIE - Centro Popolare, 18/06/2024) ribadisce che «l'autonomia differenziata non mina l'unità della Repubblica, ma riconosce le diversità che da sempre caratterizzano il nostro Paese». La parlamentare qualifica il testo come una vera e propria rivoluzione culturale capace di premiare la responsabilità e l'efficacia delle amministrazioni locali.\\[6pt] \#\# Posizioni dell'Opposizione\\[6pt] Per il Partito Democratico - Italia Democratica e Progressista, \textbf{Guerra} (Partito Democratico - Italia Democratica e Progressista, 22/12/2022) denuncia la pericolosità della manovra per l'unità del Paese dichiarando che «per potere parlare di autonomia differenziata si è finalmente capito che prima occorre definire i livelli essenziali delle prestazioni». Il gruppo critica duramente l'assenza di finanziamenti strutturali necessari per garantire concretamente i diritti di cittadinanza prima di procedere al trasferimento delle competenze.\\[6pt] Per il Movimento 5 Stelle, \textbf{Colucci} (Movimento 5 Stelle, 12/06/2024) propone di mitigare la destrutturazione del sistema per inserire «con una legge ordinaria dello Stato, un quadro unitario che possa raccordare tra di loro le varie forme di autonomia». Il partito esprime una forte preoccupazione per la frammentazione istituzionale, ritenendo che il disegno di legge trasformerà il concetto di differenziazione in una vera e propria discriminazione incostituzionale.\\[6pt] Per Alleanza Verdi e Sinistra, \textbf{Mari} (Alleanza Verdi e Sinistra, 29/04/2024) contesta apertamente la direzione ideologica del Governo chiarendo che «con la flat tax e l'autonomia differenziata dipenderà tutto solamente da dove nasci: una doppia visione, per censo e per territorio, la fine del diritto di uguaglianza». L'esponente sottolinea che un regionalismo competitivo favorirà esclusivamente il Nord, spezzando irrimediabilmente la coesione e l'unità nazionale.\\[6pt] Per Azione - Popolari Europeisti Riformatori - Renew Europe, \textbf{D'Alessio} (Azione - Popolari Europeisti Riformatori - Renew Europe, 10/07/2024) sostiene che l'approssimazione normativa del provvedimento «creerà sperequazioni, disuguaglianza, disparità di trattamento, ma anche problemi e difficoltà ai cittadini e agli imprenditori e non solo a quelli del Sud». Viene aspramente criticato l'approccio elettorale che sacrifica la semplificazione amministrativa in favore di un frettoloso accordo politico di maggioranza.\\[6pt] Per Italia Viva - Il Centro - Renew Europe, \textbf{Faraone} (Italia Viva - Il Centro - Renew Europe, 10/01/2024) osserva l'incoerenza procedurale in atto e avverte che il testo «rischia di essere esclusivamente - è anche un'opportunità da molti punti di vista - uno spot elettorale, visto che la precondizione è la realizzazione dei livelli essenziali di prestazione». Il gruppo giudica insostenibile e inattuabile il progetto senza un preventivo e trasparente stanziamento di risorse economiche.\\[6pt] Nei documenti disponibili non emergono interventi specifici del Misto sul tema. & \\
\midrule
Bioetica & \#\# Introduzione\\[6pt] Il dibattito parlamentare sui temi legati alla bioetica si è sviluppato attraverso l'esame di specifici provvedimenti, tra cui la proposta di legge A.C. 887-A sulla perseguibilità della surrogazione di maternità all'estero e il disegno di legge S. 651 in materia di alimenti coltivati in laboratorio. Le discussioni documentate hanno coinvolto oltre quindici deputati nel periodo compreso tra aprile 2023 e novembre 2023, interessando primariamente le sedute numero 85, 121, 147 e 189 della Camera.\\[6pt] \#\# Posizione del Governo\\[6pt] Per il Governo, il Ministro della Salute \textbf{Schillaci} definisce la cornice istituzionale sul fine vita e sulle cure palliative, inquadrando la soppressione del dolore come una priorità etica essenziale. \textbf{Schillaci} (Misto, 12/04/2023) afferma che occorre «garantire il diritto del malato ad accedere tanto alle cure palliative, quanto alla terapia del dolore, per assicurare il rispetto della dignità e dell'autonomia della persona umana». \\[6pt] \#\# Posizioni della Maggioranza\\[6pt] Per Fratelli d'Italia, \textbf{Cerreto} critica lo sviluppo delle tecnologie alimentari basate su tessuti cellulari artificiali, sollevando timori diretti per la sicurezza dei cittadini. \textbf{Cerreto} (Fratelli d'Italia, 06/11/2023) denuncia che «non abbiamo assolutamente alcuna certezza che queste cellule, prese da un feto, messe in un bioreattore e fatte crescere da un bioreattore, non siano dannose per la salute umana».\\[6pt] Per la Lega - Salvini Premier, \textbf{Loizzo} si oppone fermamente alla maternità surrogata, condannandola come una pratica che mercifica e sfrutta la funzione riproduttiva femminile. \textbf{Loizzo} (Lega - Salvini Premier, 26/07/2023) rileva che «l'articolo 3 della Carta dei diritti umani dell'Unione europea vieta di fare del corpo e delle sue parti una fonte di lucro».\\[6pt] Per Forza Italia - Berlusconi Presidente - PPE, \textbf{Patriarca} si esprime a favore della criminalizzazione della gestazione per altri, considerandola inaccettabile dal punto di vista bioetico e morale. \textbf{Patriarca} (Forza Italia - Berlusconi Presidente - PPE, 26/07/2023) ribadisce che «non possiamo che ritenere la maternità surrogata incompatibile con il nostro ordinamento, in quanto manifestamente contraria alla tutela della dignità umana».  Nei documenti disponibili non emergono interventi specifici di Noi Moderati (Noi con l'Italia, Coraggio Italia, UDC e Italia al Centro) - MAIE - Centro Popolare sul tema. \\[6pt] \#\# Posizioni dell'Opposizione\\[6pt] Per il Partito Democratico - Italia Democratica e Progressista, \textbf{Braga} attacca le restrizioni penali introdotte sulla surrogazione di maternità, inquadrandole come un provvedimento puramente ideologico. \textbf{Braga} (Partito Democratico - Italia Democratica e Progressista, 26/07/2023) contesta che «la vostra legge è un'arma contro tutte le famiglie omogenitoriali, contro tutti i modelli che si discostano da quello che voi considerate normale».\\[6pt] Per il Movimento 5 Stelle, \textbf{Quartini} esprime allarme per i potenziali danni collaterali che le nuove norme punitive potrebbero causare alle pratiche procreative già consentite. \textbf{Quartini} (Movimento 5 Stelle, 19/06/2023) evidenzia che «la stessa dichiarazione di illegittimità di divieto di procreazione medicalmente assistita di tipo eterologo rischia di essere vanificata».\\[6pt] Per Alleanza Verdi e Sinistra, \textbf{Grimaldi} difende l'accesso alle tecniche riproduttive come strumento di libertà, respingendo l'accusa di mercificazione del corpo sollevata dalla maggioranza. \textbf{Grimaldi} (Alleanza Verdi e Sinistra, 19/06/2023) propone la tesi secondo cui «il controllo riproduttivo non sia qualcosa di disumanizzante, che per forza mercifica la maternità e la vita, ma un'opportunità di autorealizzazione e di benessere per tutti».\\[6pt] Per Azione - Popolari Europeisti Riformatori - Renew Europe, \textbf{Pastorella} difende la validità degli studi sulla carne coltivata, invitando a riflettere oggettivamente sull'impatto etico e ambientale della nutrizione. \textbf{Pastorella} (Azione - Popolari Europeisti Riformatori - Renew Europe, 06/11/2023) sostiene sia necessario chiedersi se «la cosiddetta carne coltivata davvero mina la biodiversità o, forse, sono gli allevamenti intensivi a minare di più la biodiversità che qualcosa prodotto in laboratorio».  Nei documenti disponibili non emergono interventi specifici di Italia Viva - Il Centro - Renew Europe sul tema.\\[6pt] Per il gruppo Misto, \textbf{Magi} si batte per la legalizzazione di una forma di gestazione per altri regolamentata e altruistica, in contrapposizione ai divieti assoluti. \textbf{Magi} (Misto, 19/06/2023) auspica di «prevedere che, prima della stipula dell'accordo, vi siano verifiche della genuina volontà della donna di procedere alla gravidanza solidale». & \\
\midrule
Bolkestein & \#\# Introduzione\\[6pt] Il dibattito parlamentare sulla direttiva in questione ha interessato l'esame di svariati provvedimenti, tra cui spiccano il decreto-legge 'Milleproroghe' e il decreto Salva infrazioni'. Le discussioni hanno coinvolto oltre 20 deputati e membri dell'esecutivo in un arco temporale compreso tra il 2022 e il 2025. Le posizioni espresse nei documenti si distribuiscono lungo numerose sedute dell'Assemblea, tra cui le numero 56, 57, 193, 220, 373, 374 e 468.\\[6pt] \#\# Posizione del Governo\\[6pt] Per il Governo, il Ministro delle Infrastrutture e dei trasporti \textbf{Salvini} auspica che il settore balneare superi le lunghe incertezze legate all'interlocuzione comunitaria per giungere a una definitiva stabilità normativa. Esprimendo l'urgenza di tutelare lo sviluppo degli operatori, \textbf{Salvini} (Governo, 16/04/2025) richiede l'adozione di «un quadro di regole chiaro e sostenibile a beneficio di tutta l'Italia».\\[6pt] \#\# Posizioni della Maggioranza\\[6pt] Per Fratelli d'Italia, \textbf{Filini} rileva che l'approccio dell'esecutivo ha saggiamente protetto gli investimenti delle imprese balneari prolungando in maniera legittima i tempi prima delle procedure di gara pubbliche. Supportando l'accordo siglato con Bruxelles, \textbf{Filini} (Fratelli d'Italia, 30/10/2024) difende l'estensione «fino al 30 settembre 2027, della validità delle concessioni, prevedendo altresì la possibilità di un'ulteriore estensione».\\[6pt] Per la Lega - Salvini Premier, \textbf{Barabotti} sostiene che le prescrizioni europee in materia costituiscano un serio pericolo per il delicato tessuto produttivo composto da migliaia di microimprese familiari. Difendendo con forza i provvedimenti di rinvio adottati dal governo, \textbf{Barabotti} (Lega - Salvini Premier, 18/12/2023) definisce le aperture forzate al mercato volute dalla Bolkestein come parte di «tentativi di aggressione nei confronti della nostra economia».\\[6pt] Per Forza Italia - Berlusconi Presidente - PPE, \textbf{Bergamini} afferma che la direttiva comunitaria debba disciplinare la libera prestazione dei servizi e non la concessione di beni fisici come le spiagge. Rifiutando l'automatismo delle gare senza una previa verifica sulla reale disponibilità degli spazi costieri, \textbf{Bergamini} (Forza Italia - Berlusconi Presidente - PPE, 22/02/2023) giudica impossibile procedere senza prima accertare «una scarsità di risorse per cui bisogna liberalizzare».\\[6pt] Per Noi Moderati (Noi con l'Italia, Coraggio Italia, UDC e Italia al Centro) - MAIE - Centro Popolare, \textbf{Cavo} fa notare che le misure portate avanti dall'esecutivo intervengono efficacemente a sostegno del turismo nazionale e della continuità aziendale. Esprimendo appoggio per le tempistiche previste dai decreti, \textbf{Cavo} (Noi Moderati (Noi con l'Italia, Coraggio Italia, UDC e Italia al Centro) - MAIE - Centro Popolare, 23/02/2023) accoglie con pieno favore «diverse misure importanti legate al turismo, come la proroga delle concessioni balneari».\\[6pt] \#\# Posizioni dell'Opposizione\\[6pt] Per il Partito Democratico - Italia Democratica e Progressista, \textbf{Gnassi} denuncia l'incapacità dell'esecutivo di applicare la normativa comunitaria e di adeguare canoni giudicati insostenibilmente bassi. Criticando l'assenza di regole statali atte a limitare le concentrazioni e a favorire le piccole e medie imprese, \textbf{Gnassi} (Partito Democratico - Italia Democratica e Progressista, 29/12/2023) accusa il governo di mantenere «in una condizione di illegalità un pezzo intero di Paese».\\[6pt] Per il Movimento 5 Stelle, \textbf{Caso} contesta le proroghe indiscriminate operate dalla maggioranza, ritenendole uno strumento di privatizzazione silenziosa a netto danno del bene pubblico. Sottolineando le inequivocabili bocciature arrivate dagli organi giurisdizionali, \textbf{Caso} (Movimento 5 Stelle, 29/10/2024) ricorda a tal proposito che «le sentenze della Corte di giustizia europea e quelle del Consiglio di Stato» impongono lo svolgimento immediato di procedure selettive.\\[6pt] Per Alleanza Verdi e Sinistra, \textbf{Zanella} ribadisce la stringente necessità di fermare il rinnovo automatico delle occupazioni demaniali al fine di evitare imminenti sanzioni europee e tutelare l'accesso gratuito al mare. Deplorando l'atteggiamento governativo a strenua difesa dei privilegi acquisiti a scapito dell'erario, \textbf{Zanella} (Alleanza Verdi e Sinistra, 23/02/2023) definisce la mancata messa a bando «un grande scandalo, perché, a fronte di migliaia di chilometri di coste, per dette concessioni lo Stato incassa solo 100 milioni».\\[6pt] Per Azione - Popolari Europeisti Riformatori - Renew Europe, \textbf{Sottanelli} evidenzia la profonda incoerenza dei partiti di governo che, dopo innumerevoli promesse elettorali sull'annullamento della direttiva, sono oggi costretti a predisporne l'applicazione. Biasimando le false speranze fornite sistematicamente agli operatori del settore e l'assenza di tutele reali per la transizione, \textbf{Sottanelli} (Azione - Popolari Europeisti Riformatori - Renew Europe, 30/10/2024) ricorda criticamente come la destra avesse sempre garantito che «la Bolkestein non sarebbe stata applicata».\\[6pt] Per Italia Viva - Il Centro - Renew Europe, \textbf{Del Barba} propone di procedere celermente con l'apertura del settore marittimo per garantire la conformità dell'Italia alle basilari regole della libera concorrenza. Rifiutando la continua difesa corporativa delle rendite di posizione e richiamando i moniti del Capo dello Stato, \textbf{Del Barba} (Italia Viva - Il Centro - Renew Europe, 24/01/2024) invita categoricamente la maggioranza parlamentare a «mettere a gara le concessioni balneari».\\[6pt] Per il Misto, \textbf{Della Vedova} osserva che il mantenimento dell'attuale sistema concessorio danneggia gravemente sia le casse dello Stato sia i contribuenti, bloccando lo sviluppo di una normale competizione di mercato. Puntando il dito contro i canoni di affitto quasi irrisori incassati a fronte di enormi profitti privati, \textbf{Della Vedova} (Misto, 19/12/2023) giudica politicamente inaccettabile la volontà dell'esecutivo di «proteggere una categoria di privilegiati che vivono di rendita». & \\
\midrule
Educazione affettiva e sessuale nelle scuole & \#\# Introduzione\\[6pt] Il dibattito parlamentare sull'educazione affettiva e sessuale nelle scuole si è sviluppato tra il 2022 e il 2025, coinvolgendo oltre 40 deputati in circa 50 sedute, tra cui le sedute 185, 562, 564 e 577. La discussione si è concentrata prevalentemente sul disegno di legge Disposizioni in materia di consenso informato in ambito scolastico (A.C. 2423-A) e sulle proposte normative per il contrasto della violenza domestica e di genere.\\[6pt] \#\# Posizione del Governo\\[6pt] Per il Governo, il Ministro dell'Istruzione e del merito Giuseppe Valditara sostiene l'inclusione dell'educazione alle relazioni nei programmi vigenti di educazione civica, respingendo le accuse di censura. Valditara (Governo, 12/11/2025) evidenzia che «per la prima volta abbiamo previsto come obiettivi obbligatori di apprendimento l'educare alle relazioni, alle relazioni corrette, educare al rispetto, in particolare al rispetto verso la donna».\\[6pt] \#\# Posizioni della Maggioranza\\[6pt] Per Fratelli d'Italia, Marco Perissa sostiene che il consenso preventivo tuteli le scelte educative delle famiglie e respinge il legame diretto tra corsi scolastici e riduzione delle violenze. Perissa (Fratelli d'Italia, 10/11/2025) contesta le tesi avverse precisando che «non c'è nessuna relazione statistica fra l'educazione sessuale nelle scuole e la diminuzione dei reati di femminicidio».\\[6pt] Per la Lega - Salvini Premier, Rossano Sasso sostiene la necessità di arginare derive esterne affidando l'insegnamento esclusivamente ai docenti curricolari qualificati. Sasso (Lega - Salvini Premier, 03/12/2025) ribadisce che «con il consenso informato noi non vietiamo l'educazione affettiva, che è già prevista nelle indicazioni nazionali».\\[6pt] Per Forza Italia - Berlusconi Presidente - PPE, Rosaria Tassinari sostiene l'obbligo di consenso preventivo per bilanciare la funzione formativa scolastica con i diritti genitoriali sui temi intimi. Tassinari (Forza Italia - Berlusconi Presidente - PPE, 03/12/2025) rileva che «si tratta di riconoscere il diritto e la responsabilità delle famiglie di partecipare pienamente alle scelte educative che incidono su aspetti tanto delicati della formazione dei figli».\\[6pt] Per Noi Moderati (Noi con l'Italia, Coraggio Italia, UDC e Italia al Centro) - MAIE - Centro Popolare, Maurizio Enzo Lupi sostiene il primato educativo costituzionale della famiglia, opponendosi a una totale delega dello Stato alla scuola. Lupi (Noi Moderati (Noi con l'Italia, Coraggio Italia, UDC e Italia al Centro) - MAIE - Centro Popolare, 03/12/2025) osserva che «nessuno può togliere questa potestà alla famiglia e ai genitori».\\[6pt] \#\# Posizioni dell'Opposizione\\[6pt] Per il Partito Democratico - Italia Democratica e Progressista, Valentina Ghio sostiene che ostacolare i progetti scolastici attraverso la burocrazia del consenso indebolisca il contrasto strutturale alla violenza di genere. Ghio (Partito Democratico - Italia Democratica e Progressista, 11/11/2025) denuncia che «non garantire un accesso equo e uniforme ai percorsi educativi su affettività, relazioni e sessualità significa proprio indebolire la capacità dello Stato di prevenire comportamenti violenti e discriminatori».\\[6pt] Per il Movimento 5 Stelle, Stefania Ascari sostiene l'introduzione obbligatoria della materia fin dalla scuola primaria per scardinare precocemente gli stereotipi patriarcali. Ascari (Movimento 5 Stelle, 23/11/2022) propone di «introdurre da subito in modo sistemico e continuativo l'insegnamento dell'educazione affettiva e sessuale dai primi banchi di scuola».\\[6pt] Per Alleanza Verdi e Sinistra, Elisabetta Piccolotti sostiene che i dati scientifici internazionali certifichino l'efficacia dei corsi strutturali nel diminuire abusi e malattie trasmissibili. Piccolotti (Alleanza Verdi e Sinistra, 10/11/2025) sottolinea che gli studi medici «mostrano risultati inequivocabili dell'educazione sessuale strutturale a scuola».\\[6pt] Per Azione - Popolari Europeisti Riformatori - Renew Europe, Elena Bonetti sostiene che il requisito del consenso informato paralizzi l'autonomia didattica rendendo facoltative importanti tematiche di salute pubblica. Bonetti (Azione - Popolari Europeisti Riformatori - Renew Europe, 03/12/2025) afferma che «quando questo diventerà norma primaria, sarà prevalente anche rispetto a tutte le attività curriculari».\\[6pt] Per Italia Viva - Il Centro - Renew Europe, Roberto Giachetti sostiene che l'assenza di un chiaro indirizzo politico nazionale stia generando disuguaglianze inaccettabili e sistemiche tra gli istituti. Giachetti (Italia Viva - Il Centro - Renew Europe, 11/11/2025) auspica un intervento normativo ricordando che «il rapporto tra la scuola italiana e l'educazione sessuale è fermo da 50 anni».\\[6pt] Per il Misto, Riccardo Magi sostiene che l'educazione sessuo-affettiva pubblica assicuri tutele ed emancipazione essenziali agli adolescenti privi di un supporto familiare adeguato. Magi (Misto, 03/12/2025) dichiara che «proprio per questo esiste la scuola, per garantire pari opportunità anche nell'educazione affettiva e sessuale, con linguaggi corretti». & \\
\midrule
Gestione dei Flussi Migratori & \#\# Introduzione\\[6pt] L'analisi si basa sui resoconti stenografici della Camera dei Deputati relativi a molteplici provvedimenti, tra cui i disegni di legge di conversione dei decreti-legge n. 1/2023, n. 20/2023 e n. 37/2025. Il dibattito ha coinvolto oltre sessanta deputati intervenuti in centinaia di sedute parlamentari, tra cui le sedute n. 48, n. 95 e n. 554, svoltesi nel periodo compreso tra novembre 2022 e dicembre 2025. I documenti esaminati coprono sia le discussioni sulle informative urgenti del Governo, sia la ratifica di protocolli internazionali in materia di asilo e controllo delle frontiere. Posizione del Governo Per il Governo, Giorgia Meloni chiarisce che la gestione delle frontiere meridionali richiede una forte collaborazione comunitaria e partnership strategiche con le nazioni africane per fermare tempestivamente i trafficanti. La Presidente del Consiglio (Fratelli d'Italia, 13/12/2022) afferma che «Il fianco Sud della sfida migratoria non è meno importante del fianco Est».\\[6pt] \#\# Posizioni della Maggioranza\\[6pt] Per Fratelli d'Italia, Sara Kelany ribadisce che le priorità dell'esecutivo si fondano sul rigido controllo degli ingressi e sugli accordi per trasferire le procedure di accoglienza al di fuori dei confini nazionali. La deputata (Fratelli d'Italia, 19/03/2025) afferma che le scelte del Governo convergono nella «lotta senza quartiere ai trafficanti di esseri umani; difesa delle frontiere; esternalizzazione della gestione dei flussi verso Paesi terzi».\\[6pt] Per la Lega - Salvini Premier, Igor Giancarlo Iezzi denuncia il ricorso abusivo agli strumenti di tutela umanitaria, inquadrandoli come meccanismi che incentivano la permanenza irregolare di chi non possiede i requisiti per l'asilo. Il legislatore (Lega - Salvini Premier, 04/05/2023) afferma che «La protezione speciale allargata, invece, è un regalo del Partito Democratico, una vera e propria sanatoria, utile solo a regalare un permesso a chi non ha nessun diritto di stare qui».\\[6pt] Per Forza Italia - Berlusconi Presidente - PPE, Paolo Emilio Russo auspica un sistema normativo incentrato sulle espulsioni rapide per chi delinque e sull'apertura parallela di canali per i lavoratori stranieri regolari. Il parlamentare (Forza Italia - Berlusconi Presidente - PPE, 16/10/2024) afferma che «La corretta gestione dei flussi migratori, consentendo l'immigrazione ai lavoratori regolari, riconoscendo asilo agli aventi diritto e respingendo, tramite rimpatri celeri, gli irregolari è di fondamentale importanza».\\[6pt] Per Noi Moderati (Noi con l'Italia, Coraggio Italia, UDC e Italia al Centro) - MAIE - Centro Popolare, Alessandro Colucci propone di considerare le iniziative di solidarietà per i veri profughi come strettamente complementari alla ferma opposizione verso gli sbarchi illegali. Il deputato (Noi Moderati (Noi con l'Italia, Coraggio Italia, UDC e Italia al Centro) - MAIE - Centro Popolare, 28/11/2023) afferma che è prioritario «accogliere i migranti che scappano da guerre, da persecuzioni, da fame, e contrastare l'immigrazione clandestina».\\[6pt] \#\# Posizioni dell'Opposizione\\[6pt] Per il Partito Democratico - Italia Democratica e Progressista, Laura Boldrini contesta le recenti normative governative poiché smantellano i modelli di inclusione locale per sostituirli con strutture detentive inadeguate. La deputata (Partito Democratico - Italia Democratica e Progressista, 04/05/2023) afferma che «I richiedenti asilo non potranno essere più inseriti nei circuiti della rete SAI - costituita da piccoli centri che favoriscono l'integrazione, anche perché non hanno l'impatto pesante, incisivo e invasivo sul territorio - e potranno, invece, essere trattenuti nei centri di detenzione, negli hotspot e addirittura nei CPR».\\[6pt] Per il Movimento 5 Stelle, Raffaele Bruno rileva la stringente necessità di coinvolgere le istituzioni continentali per superare le risposte puramente emergenziali e pianificare ingressi sicuri. Il parlamentare (Movimento 5 Stelle, 29/11/2023) afferma che risulta essenziale «promuovere canali d'ingresso legali attraverso una progressiva programmazione di flussi di lavoratori a livello europeo».\\[6pt] Per Alleanza Verdi e Sinistra, Devis Dori evidenzia il totale dissenso verso le intese internazionali di contenimento che violano i diritti fondamentali, esigendo un approccio strutturale incentrato sulla tutela della vita. L'onorevole (Alleanza Verdi e Sinistra, 02/05/2023) afferma che «Servono corridoi umanitari, garanzia dell'accesso alla procedura di asilo e all'accoglienza, abbandono delle politiche di esternalizzazione e dei loro scellerati risultati».\\[6pt] Per Azione - Popolari Europeisti Riformatori - Renew Europe, Giulia Pastorella argomenta che le procedure basate sulle quote occasionali siano inadeguate a soddisfare il reale fabbisogno del mercato del lavoro italiano. La deputata (Azione - Popolari Europeisti Riformatori - Renew Europe, 14/05/2025) afferma che «È necessaria una gestione seria e pragmatica di questi flussi migratori, ampliando i canali regolari, non basta il solito decreto flussi con il click day».\\[6pt] Per Italia Viva - Il Centro - Renew Europe, Roberto Giachetti osserva come i provvedimenti dell'esecutivo contengano gravi forzature lesive della privacy e delle garanzie costituzionali per i cittadini stranieri sottoposti a controlli. L'esponente (Italia Viva - Il Centro - Renew Europe, 26/11/2024) afferma che «questo testo, con l'articolo 12, permette di accedere immediatamente ai dispositivi elettronici».\\[6pt] Per il gruppo Misto, Dieter Steger precisa che le espulsioni dei soggetti non aventi diritto sono condivisibili solo se integrate in un coerente sistema europeo di governo degli ingressi e di cooperazione tra Stati. Il legislatore (Misto, 15/05/2025) afferma che «Comprendiamo e condividiamo l'esigenza di rafforzare i rimpatri, di contrastare gli abusi di diritto d'asilo, di impedire che l'illegalità diventi la scorciatoia per entrare in Europa». & \\
\midrule
Legge di Bilancio & \#\# Introduzione\\[6pt] Il documento analizza il dibattito parlamentare relativo alla Legge di Bilancio e ai collegati documenti di programmazione economica, svoltosi tra il mese di novembre 2022 e la fine del 2025. Nel periodo in esame, che include sedute specifiche come la numero 27, la 402 e la 588, si registrano decine di interventi da parte di oltre venti deputati appartenenti a tutti gli schieramenti politici. I lavori hanno coinvolto attivamente sia i rappresentanti dell'Esecutivo sia i vari gruppi parlamentari in molteplici discussioni sulle coperture e sui saldi di finanza pubblica.\\[6pt] Posizione del Governo\\[6pt] Per il Governo, il Ministro dell'Economia e delle finanze Giorgetti incentra la sua visione sull'espansione controllata e sul rigoroso monitoraggio dei conti pubblici. A tal riguardo, Giorgetti (Lega - Salvini Premier, 13/11/2024) dichiara che «nell'anno 2024 l'andamento della spesa ha mostrato una curva progressivamente crescente, che, se sarà confermata, come prevediamo e auspichiamo, consentirà di raggiungere un livello di spesa superiore ai 20 miliardi, coerente con l'ultima stima di finanza pubblica».\\[6pt] Posizioni della Maggioranza\\[6pt] Per Fratelli d'Italia, l'assetto contabile della manovra è considerato essenziale per risanare le finanze e tutelare simultaneamente i cittadini. In merito, Giorgianni (Fratelli d'Italia, 04/12/2024) rileva che «Mettere in ordine i conti pubblici, ridurre le spese improduttive per rilanciare la produzione, incentivare gli investimenti e sostenere soprattutto il potere d'acquisto delle famiglie: questi sono gli obiettivi chiave di questa manovra».\\[6pt] Per la Lega - Salvini Premier, le scelte programmatiche attuali si pongono in perfetta coerenza con un percorso triennale di tutela economica delle fasce vulnerabili. Su questo punto, Barabotti (Lega - Salvini Premier, 28/12/2025) evidenzia che «Questa legge di bilancio, difatti, non nasce dal nulla, ma si inserisce in un percorso avviato fin dall'inizio della legislatura fatto di tre manovre che consecutivamente hanno messo al centro le famiglie più fragili».\\[6pt] Per Forza Italia - Berlusconi Presidente - PPE, le misure necessitano di una profonda e rapida azione di revisione della spesa per garantire un utilizzo ottimale delle risorse. Di conseguenza, Mangialavori (Forza Italia - Berlusconi Presidente - PPE, 27/04/2023) auspica che «questa opera possa concludersi in tempi brevi, in modo che i suoi esiti siano condivisi quanto prima».\\[6pt] Per Noi Moderati (Noi con l'Italia, Coraggio Italia, UDC e Italia al Centro) - MAIE - Centro Popolare, il testo in discussione rappresenta un bilanciamento ideale tra rigore strutturale e supporto al tessuto sociale. Difatti, Cavo (Noi Moderati (Noi con l'Italia, Coraggio Italia, UDC e Italia al Centro) - MAIE - Centro Popolare, 28/12/2025) ribadisce che «Il Governo ha saputo presentare una legge di bilancio che coniuga rigorosa prudenza nei conti pubblici con interventi mirati e concreti a sostegno di famiglie, lavoratori, imprese e welfare».\\[6pt] Posizioni dell'Opposizione\\[6pt] Per il Partito Democratico - Italia Democratica e Progressista, le coperture della manovra sono ottenute a discapito di settori cruciali e attraverso rinunce strategiche sulle entrate. Per queste ragioni, Braga (Partito Democratico - Italia Democratica e Progressista, 23/12/2022) contesta che «In questa legge di bilancio voi avete scelto di finanziare alcune cose - vedi gli 800 milioni per le società di calcio - di definanziarne altre, di rinunciare alle entrate della lotta all'evasione fiscale».\\[6pt] Per il Movimento 5 Stelle, il testo programmatico trascura deliberatamente settori fondamentali per il rilancio del Paese come le retribuzioni. Entrando nel merito, Dell'Olio (Movimento 5 Stelle, 29/12/2025) denuncia che «Nel Paese ci sono quattro grandi emergenze che questa legge di bilancio ignora completamente. La prima: i salari».\\[6pt] Per Alleanza Verdi e Sinistra, le decisioni governative colpiscono in modo sproporzionato i dicasteri legati allo sviluppo infrastrutturale e alla mobilità sostenibile. Analizzando i fondi, Grimaldi (Alleanza Verdi e Sinistra, 28/12/2025) sottolinea che «Il MIT è il Ministero più penalizzato. Nel 2026 le infrastrutture perdono 524 milioni di euro».\\[6pt] Per Azione - Popolari Europeisti Riformatori - Renew Europe, il documento risulta destrutturato e deficitario di un orizzonte organico in grado di favorire le riforme. A tal proposito, Grippo (Azione - Popolari Europeisti Riformatori - Renew Europe, 19/12/2024) afferma che «ci troviamo a discutere una legge di bilancio ampiamente riscritta, con una pioggia di emendamenti senza relazioni tecniche e copertura».\\[6pt] Per Italia Viva - Il Centro - Renew Europe, la manovra si caratterizza per l'assenza di priorità chiare e per un approccio prettamente contabile senza spinta espansiva. In sintesi, Del Barba (Italia Viva - Il Centro - Renew Europe, 20/12/2024) osserva che «ci apprestiamo a votare una legge di bilancio scialba, priva di direzione o anche solo di singoli spunti».\\[6pt] Per il Gruppo Misto, l'esigenza di liberare risorse per la crescita deve tradursi in una valutazione rigorosa e pubblica della qualità della spesa infrastrutturale. Con questo spirito, Steger (Misto, 06/08/2025) propone che «ogni progetto sia accompagnato da analisi costi-benefici pubblica e pubblicata in un portale open data perché la trasparenza non è un orpello, è garanzia di fiducia» & \\
\midrule
Politiche sull'Intelligenza Artificiale & \#\# Introduzione\\[6pt] Il dibattito parlamentare analizzato si è concentrato prevalentemente sul disegno di legge S. 1146 recante disposizioni e deleghe al Governo in materia di intelligenza artificiale (A.C. 2316-A) e sui decreti connessi al potenziamento della pubblica amministrazione. Le discussioni si sono svolte tra il 2023 e il 2026, coinvolgendo in modo trasversale oltre quaranta deputati nel corso di numerose sedute, con particolare intensità durante le sedute 497, 498 e 499. L'impianto documentale raccoglie una vasta mole di interventi focalizzati sugli impatti etici, occupazionali, normativi e industriali legati allo sviluppo tecnologico.\\[6pt] Posizioni della Maggioranza\\[6pt] Per Fratelli d'Italia, Russo sostiene la necessità di bilanciare le immense opportunità economiche della tecnologia con norme rigorose volte a scongiurare usi impropri o lesivi dei diritti fondamentali. Zucconi (Fratelli d'Italia, 25/06/2025) evidenzia che «il tema centrale dell'uso dell'intelligenza artificiale si basa sul tema della contraffazione e della manipolazione della realtà».\\[6pt] Per Lega - Salvini Premier, Gusmeroli reputa fondamentale governare la transizione digitale proteggendo il tessuto produttivo delle piccole e medie imprese e favorendo l'inclusione dei cittadini meno alfabetizzati. Barabotti (Lega - Salvini Premier, 25/06/2025) auspica che «il nostro compito come legislatori è fare in modo che questo impianto normativo non resti sulla carta».\\[6pt] Per Forza Italia - Berlusconi Presidente - PPE, Tenerini rifiuta l'idea di un controllo cieco sui lavoratori, difendendo invece il principio antropocentrico di uno sviluppo tecnologico basato sulla protezione intelligente della persona. Caroppo (Forza Italia - Berlusconi Presidente - PPE, 25/06/2025) ribadisce che «l'intelligenza artificiale non sostituirà le persone, ma le affiancherà e le potenzierà».\\[6pt] Per Noi Moderati (Noi con l'Italia, Coraggio Italia, UDC e Italia al Centro) - MAIE - Centro Popolare, Colucci richiede un rapido aggiornamento delle competenze dei dipendenti pubblici per affrontare l'innovazione senza subire passivamente l'automazione. Cavo (Noi Moderati (Noi con l'Italia, Coraggio Italia, UDC e Italia al Centro) - MAIE - Centro Popolare, 25/06/2025) propone che «questo provvedimento rappresenta un passo cruciale verso la protezione dei dati personali, con particolare attenzione, appunto, all'accesso dei minori, all'uso in ambito sanitario, alla ricerca scientifica e al lavoro».\\[6pt] Posizioni dell'Opposizione\\[6pt] Per il Partito Democratico - Italia Democratica e Progressista, Casu invoca l'istituzione di un'autorità terza indipendente per la gestione del settore, opponendosi fermamente all'affidamento delle competenze ad agenzie governative ritenute non autonome. Roggiani (Partito Democratico - Italia Democratica e Progressista, 24/06/2025) denuncia che «aziende, magari straniere, possano in modo illegale saccheggiare le opere dei nostri autori e trarne un profitto che non è generato dalla loro creatività».\\[6pt] Per il Movimento 5 Stelle, L'Abbate suggerisce l'introduzione di certificazioni obbligatorie per i fornitori nei contratti pubblici al fine di impedire che algoritmi non verificati prendano decisioni su cittadini e imprese. Ferrara (Movimento 5 Stelle, 23/06/2025) contesta che «state ignorando la necessità di prevedere investimenti adeguati a formare il personale».\\[6pt] Per Alleanza Verdi e Sinistra, Piccolotti esige il pieno controllo statale sui dati sensibili della pubblica amministrazione, osteggiando categoricamente la conservazione delle informazioni su server esteri o l'accesso incontrollato da parte dei privati. Ghirra (Alleanza Verdi e Sinistra, 24/06/2025) rileva che «i dati della pubblica amministrazione dovrebbero essere, a nostro avviso, considerati un bene di interesse pubblico in modo da disciplinarne un corretto e legittimo utilizzo».\\[6pt] Per Azione - Popolari Europeisti Riformatori - Renew Europe, Pastorella dubita dell'effettiva capacità delle amministrazioni locali di gestire gli adempimenti in materia cibernetica previsti senza un chiaro quadro di competenze professionali strutturate. Benzoni (Azione - Popolari Europeisti Riformatori - Renew Europe, 25/06/2025) afferma che «è un provvedimento che ha raccolto molte delle nostre osservazioni, ma che poco agisce, nella maniera pratica, su quello che è il grande tema del futuro».\\[6pt] Per Italia Viva - Il Centro - Renew Europe, Giachetti critica l'inefficienza delle politiche attuali sull'ammodernamento digitale, ritenendole in potenziale disaccordo con la costruzione di una pubblica amministrazione moderna e realmente competitiva. Del Barba (Italia Viva - Il Centro - Renew Europe, 25/06/2025) rimarca che «i dati sono l'oro dell'intelligenza artificiale».\\[6pt] Per il gruppo Misto, Calderone giudica prioritario puntare sulla formazione dei lavoratori in condizione di vulnerabilità, integrando l'acquisizione di nuove competenze digitali come argine contro l'esclusione sociale. Steger (Misto, 25/06/2025) chiarisce che «l'intelligenza artificiale, se governata con responsabilità e con visione, non è una minaccia, ma può diventare la grande alleata della democrazia». & \\
\midrule
Riarmo europeo & \#\# Introduzione\\[6pt] Il dibattito parlamentare sul piano di riarmo europeo e sulla partecipazione alle missioni internazionali si è articolato in molteplici sedute tra la fine del 2022 e la fine del 2025, trovando il suo culmine nella seduta 465 del 10 aprile 2025 dedicata specificamente all'esame delle mozioni sul tema. I resoconti analizzati documentano oltre sessanta interventi che hanno coinvolto più di quaranta deputati, garantendo un'ampia rappresentazione delle diverse sensibilità istituzionali.\\[6pt] Posizione del Governo\\[6pt] Il Ministro dell'Economia e delle finanze Giancarlo Giorgetti sottolinea l'impegno istituzionale per finanziare il rafforzamento della sicurezza continentale massimizzando le sinergie e senza penalizzare i servizi ai cittadini. Giorgetti (Lega - Salvini Premier, 12/03/2025) chiarisce che la proposta italiana all'Ecofin mira a colmare il divario di investimenti «cercando, per quanto possibile, di minimizzare l'impatto sul debito pubblico».\\[6pt] Posizioni della Maggioranza\\[6pt] Per Fratelli d'Italia, Giangiacomo Calovini rileva che sostenere una politica di potenziamento della capacità difensiva non significa incoraggiare logiche belliciste, ma dotare il continente degli strumenti necessari a salvaguardare la pace. Calovini (Fratelli d'Italia, 10/04/2025) rimarca che un'Unione proiettata nel mondo «non può rimanere dipendente da altri per la propria sicurezza».\\[6pt] Per la Lega - Salvini Premier, Simone Billi contesta con grande fermezza la decisione di contrarre ulteriore indebitamento comunitario per finanziare il comparto militare, suggerendo invece di indirizzare i fondi verso ospedali e scuole. Billi (Lega - Salvini Premier, 10/04/2025) specifica in modo diretto che «ci opponiamo fermamente a questi 800 miliardi di debiti per la difesa europea».\\[6pt] Per Forza Italia - Berlusconi Presidente - PPE, Paolo Barelli auspica un salto di qualità nella gestione delle sfide globali, valutando le nuove risorse alla sicurezza come strumenti imprescindibili per tutelare le catene di approvvigionamento. Barelli (Forza Italia - Berlusconi Presidente - PPE, 19/03/2025) nota tuttavia che «difesa europea è un termine più appropriato di riarmo europeo, perché difesa significa molto altro».\\[6pt] Per Noi Moderati (Noi con l'Italia, Coraggio Italia, UDC e Italia al Centro) - MAIE - Centro Popolare, Maria Rosaria Carfagna dichiara che l'architettura geopolitica nata dopo i conflitti mondiali è esaurita e il continente deve costruire una deterrenza capace di scoraggiare mire espansionistiche. Carfagna (Noi Moderati (Noi con l'Italia, Coraggio Italia, UDC e Italia al Centro) - MAIE - Centro Popolare, 19/03/2025) ricorda inoltre che «non penso che disarmo possa essere uno slogan».\\[6pt] Posizioni dell'Opposizione\\[6pt] Per il Partito Democratico - Italia Democratica e Progressista, Elena Ethel Schlein denuncia l'architettura finanziaria dell'attuale piano von der Leyen, accusandolo di stimolare logiche nazionali sugli investimenti invece di promuovere programmi congiunti. Schlein (Partito Democratico - Italia Democratica e Progressista, 19/03/2025) avverte che il progetto così strutturato «rischia di ritardare l'obiettivo condiviso di una difesa comune».\\[6pt] Per il Movimento 5 Stelle, Arnaldo Lomuti propone di abbandonare l'inseguimento acritico dei target di spesa della NATO, concentrando gli sforzi per superare la frammentazione produttiva dell'industria militare continentale. Lomuti (Movimento 5 Stelle, 22/01/2025) valuta negativamente le duplicazioni dei mezzi e suggerisce che «basterebbe mettere in comune queste immense risorse, perché l'Europa diventasse la seconda o terza potenza militare globale».\\[6pt] Per Alleanza Verdi e Sinistra, Luana Zanella ribadisce la propria radicale contrarietà all'aumento delle spese per gli armamenti, individuando nel piano in discussione un rischio concreto di deviazione dei fondi sociali verso la logica dello scontro bellico. Zanella (Alleanza Verdi e Sinistra, 07/04/2025) esprime la profonda preoccupazione che si finisca per «favorire l'economia di guerra, il riarmo di ogni singolo Paese».\\[6pt] Per Azione - Popolari Europeisti Riformatori - Renew Europe, Federica Onori evidenzia l'estrema urgenza di rispettare gli impegni atlantici e invita le forze politiche a non sottrarsi alle responsabilità di fronte alle minacce territoriali all'Unione. Onori (Azione - Popolari Europeisti Riformatori - Renew Europe, 07/04/2025) esorta apertamente l'Esecutivo a «partecipare attivamente al ReArm Europe».\\[6pt] Per Italia Viva - Il Centro - Renew Europe, Mauro Del Barba osserva che la necessità di incrementare la capacità di deterrenza non deve mai tradursi in nuovi gravami fiscali o in tagli letali al sistema sanitario nazionale. Del Barba (Italia Viva - Il Centro - Renew Europe, 10/09/2025) accetta l'impegno strategico ma precisa che «si rifiutano soluzioni pasticciate, prive di organicità o meramente burocratiche, che allontanano il modello della difesa comune europea».\\[6pt] Per il Misto, Dieter Steger afferma che la brutale aggressione russa all'Ucraina ha distrutto ogni illusione pacifista, imponendo la creazione urgente di un mercato unico militare per affrancarsi dalle dipendenze critiche ed esterne. Steger (Misto, 10/04/2025) difende in maniera convinta le scelte comunitarie sostenendo che «il piano ReArm Europe rappresenta un'occasione storica per costruire un mercato unico europeo della difesa». & \\
\midrule
Riforma del Premierato & \#\# Introduzione\\[6pt] Il dibattito parlamentare inerente al disegno di legge sulla revisione della forma di governo, comunemente noto come Riforma del Premierato, si è sviluppato trasversalmente durante l'attuale legislatura registrando la partecipazione attiva di oltre venti deputati in Aula. Le discussioni si sono articolate in diverse sedute tematiche, tra cui le numero 115, 267, 285 e 531, svoltesi in un arco temporale compreso tra giugno 2023 e settembre 2025.\\[6pt] Posizione del Governo\\[6pt] Per il Governo, la Ministra per le Riforme istituzionali e la semplificazione normativa Alberti Casellati Maria sostiene la necessità di accrescere la stabilità dell'esecutivo tutelando parallelamente le prerogative delle Camere. Nello specifico, Alberti Casellati Maria (Forza Italia - Berlusconi Presidente - PPE, 23/04/2025) rileva che «nella cornice del premierato [...] la conservazione del modello bicamerale perfetto sia la più idonea a preservare e rafforzare quella centralità che la Costituzione riconosce al Parlamento».\\[6pt] Posizioni della Maggioranza\\[6pt] Per Fratelli d'Italia, Tremaglia supporta convintamente l'iniziativa definendola un passaggio essenziale per garantire al Paese un esecutivo capace di intervenire direttamente secondo il chiaro mandato popolare. A tal proposito, Michelotti (Fratelli d'Italia, 28/05/2024) auspica che «si vada, a grandi passi, verso la riforma di Governo, del Premierato, che sta proseguendo l'esame in Senato».\\[6pt] Per la Lega - Salvini Premier, Barabotti difende l'elezione diretta del Presidente del Consiglio ritenendola una misura fondamentale per garantire solidi orizzonti di governabilità all'intero sistema. Di conseguenza, Ziello (Lega - Salvini Premier, 29/02/2024) afferma che «Cosa c'entra il premierato, un'elezione diretta del Presidente del Consiglio con la volontà di reprimere il dissenso, che non esiste?».\\[6pt] Per Forza Italia - Berlusconi Presidente - PPE, Battilocchio supporta apertamente il disegno di legge costituzionale ricordando il forte apporto costruttivo della propria forza politica alla stesura del testo originario. In quest'ottica, Battilocchio (Forza Italia - Berlusconi Presidente - PPE, 29/04/2024) ribadisce che «al Senato sta per approdare in Aula la riforma del cosiddetto Premierato, riforme alle quali Forza Italia ha contribuito attivamente».\\[6pt] Nei documenti disponibili non emergono interventi specifici di Noi Moderati (Noi con l'Italia, Coraggio Italia, UDC e Italia al Centro) - MAIE - Centro Popolare sul tema.\\[6pt] Posizioni dell'Opposizione\\[6pt] Per il Partito Democratico - Italia Democratica e Progressista, Braga si oppone duramente al disegno di legge ritenendolo una minaccia assoluta per il mantenimento delle garanzie democratiche. Conseguentemente, Schlein (Partito Democratico - Italia Democratica e Progressista, 18/06/2024) contesta che «concentra tutti i poteri in mano al Capo del Governo».\\[6pt] Per il Movimento 5 Stelle, Baldino critica la riforma considerandola un mezzo per instaurare un sistema accentratore a discapito delle delicate funzioni di bilanciamento dello Stato. Su questo punto, Auriemma (Movimento 5 Stelle, 23/10/2024) evidenzia che «va a stravolgere completamente il nostro assetto costituzionale, incidendo in maniera negativa sul nostro garante della Repubblica».\\[6pt] Per Alleanza Verdi e Sinistra, Grimaldi respinge il progetto di legge definendolo un attacco diretto all'equilibrio dei poteri accuratamente previsto dai padri costituenti. In sintonia con questa posizione, Zanella (Alleanza Verdi e Sinistra, 16/09/2025) denuncia che «Con il premierato è chiaro il disegno di accentrare il potere nelle mani del Governo, prosciugare ulteriormente le competenze e le funzioni in capo al Parlamento».\\[6pt] Per Azione - Popolari Europeisti Riformatori - Renew Europe, Richetti censura l'approccio istituzionale accusando l'esecutivo di voler limitare le storiche funzioni legislative. Nel dettaglio, Bonetti (Azione - Popolari Europeisti Riformatori - Renew Europe, 29/04/2024) dichiara che «entrambe queste riforme, sia l'autonomia differenziata che il Premierato, vanno ad annichilire completamente un effettivo ruolo del Parlamento».\\[6pt] Per Italia Viva - Il Centro - Renew Europe, Giachetti interpreta l'intera manovra di revisione costituzionale come frutto di un mero baratto politico tra le forze della compagine governativa. Riguardo alle implicazioni tecniche, Faraone (Italia Viva - Il Centro - Renew Europe, 18/06/2024) propone che «quando si decide di bloccare una riforma che prevedeva l'elezione diretta del Premier [...] naturalmente, questo doveva prevedere una riforma elettorale che fosse coerente».\\[6pt] Nei documenti disponibili non emergono interventi specifici di Misto sul tema. & \\
\midrule
Riforma della Giustizia & \#\# Introduzione\\[6pt] I documenti in esame coprono un arco temporale dal dicembre 2022 al gennaio 2026, comprendendo numerose sedute parlamentari, tra cui le numero 39, 69, 228, 410, 531 e 598. Il dibattito ha coinvolto oltre sessanta deputati in svariati interventi incentrati su provvedimenti cardine come il disegno di legge costituzionale A.C. 1917, le comunicazioni ministeriali e i decreti-legge per l'attuazione del PNRR. Le discussioni si sono focalizzate prevalentemente sull'efficienza procedurale, sull'ordinamento penitenziario e sulla separazione delle funzioni giudiziarie.\\[6pt] Posizione del Governo\\[6pt] Per il Governo, il Ministro della Giustizia Nordio rileva i progressi fatti nella digitalizzazione e nella riduzione delle pendenze nel settore civile e penale per centrare i parametri strutturali richiesti a livello europeo. Nordio (Governo, 27/03/2024) afferma che «la separazione delle carriere è consustanziale a una riforma, peraltro, forse, ancora più importante, che è quella del Consiglio superiore della magistratura».\\[6pt] Posizioni della Maggioranza\\[6pt] Per Fratelli d'Italia, Dondi ribadisce che la revisione costituzionale mira a tutelare la totale imparzialità del giudizio e a sradicare le influenze delle correnti all'interno del sistema giudiziario. Dondi (Fratelli d'Italia, 16/09/2025) afferma che «Fratelli d'Italia porta avanti con coerenza ciò che ha sempre promesso ai cittadini, separare le carriere, rafforzare la terzietà del giudice, introdurre l'Alta Corte disciplinare».\\[6pt] Per la Lega - Salvini Premier, Bisa propone una netta distinzione tra chi indaga e chi emette sentenze al fine di responsabilizzare la pubblica accusa e garantire piene garanzie ai cittadini. Bisa (Lega - Salvini Premier, 21/01/2026) afferma che «non può esistere un processo equo, se non vi è una netta distinzione tra chi esercita l'azione penale e chi è chiamato a giudicare».\\[6pt] Per Forza Italia - Berlusconi Presidente - PPE, Patriarca auspica un rafforzamento del rito accusatorio come condizione imprescindibile per la parità tra le parti processuali e le libertà dell'imputato. Patriarca (Forza Italia - Berlusconi Presidente - PPE, 22/01/2025) afferma che «la separazione delle carriere è una garanzia per la giustizia ed è una misura finalizzata a rafforzare la terzietà del giudice e a migliorare la qualità del sistema giudiziario».\\[6pt] Per Noi Moderati (Noi con l'Italia, Coraggio Italia, UDC e Italia al Centro) - MAIE - Centro Popolare, Colucci evidenzia la necessità di interrompere la promiscuità formativa tra le diverse funzioni per restituire credibilità istituzionale e prevenire squilibri di potere. Colucci (Noi Moderati (Noi con l'Italia, Coraggio Italia, UDC e Italia al Centro) - MAIE - Centro Popolare, 16/09/2025) afferma che «separare le carriere significa depoliticizzare la giustizia e non politicizzarla».\\[6pt] Posizioni dell'Opposizione\\[6pt] Per il Partito Democratico - Italia Democratica e Progressista, Gianassi denuncia l'intento governativo di spaccare l'ordine giudiziario per colpirne l'autonomia, trascurando priorità reali come i tempi dilatati dei processi e le carenze di organico. Gianassi (Partito Democratico - Italia Democratica e Progressista, 09/12/2024) afferma che «oggi esiste già la separazione delle carriere».\\[6pt] Per il Movimento 5 Stelle, Cafiero De Raho contesta la frammentazione della giurisdizione e la duplicazione degli organi di autogoverno, considerandole azioni dirette a indebolire la magistratura di fronte all'esecutivo. Cafiero De Raho (Movimento 5 Stelle, 16/01/2025) afferma che «la separazione delle carriere determinerà l'allontanamento dei pubblici ministeri dalla cultura della giurisdizione».\\[6pt] Per Alleanza Verdi e Sinistra, Zaratti osserva il grave pericolo di trasformare la figura del magistrato requirente in un ruolo sottomesso alla politica e prossimo alle forze di polizia, smantellando il modello costituzionale. Zaratti (Alleanza Verdi e Sinistra, 09/01/2025) afferma che «dietro questo progetto di riforma vi è la palese ispirazione di abbattere il livello di indipendenza reale della magistratura».\\[6pt] Per Azione - Popolari Europeisti Riformatori - Renew Europe, D'Alessio valuta positivamente una divisione equilibrata tra funzioni giudicanti e requirenti, chiedendo però investimenti massicci nelle risorse umane e strutturali per risolvere le inefficienze degli uffici. D'Alessio (Azione - Popolari Europeisti Riformatori - Renew Europe, 17/01/2024) afferma che «i problemi si risolvono solo se la giustizia funziona realmente, con organici sufficienti e un'organizzazione degna di un Paese civile».\\[6pt] Per Italia Viva - Il Centro - Renew Europe, Boschi dichiara l'astensione del suo gruppo poiché il testo in discussione, pur affrontando un tema formalmente condiviso, risulta deformato in modo da generare conflitti tra magistrati. Boschi (Italia Viva - Il Centro - Renew Europe, 16/09/2025) afferma che «questa è una riforma di una parte della magistratura contro un'altra parte della magistratura».\\[6pt] Per il gruppo Misto, Della Vedova sottolinea l'urgenza di focalizzarsi sullo scardinamento del sistema correntizio, individuando nelle riforme della fase preliminare del processo la vera priorità per limitare gli errori e le distorsioni. Della Vedova (Misto, 16/01/2025) afferma che «l'obiettivo di disarticolare il regime correntizio dentro la magistratura è certamente un obiettivo condiviso». & \\
\midrule
Ritorno al Nucleare & \#\# Introduzione\\[6pt] Il dibattito parlamentare sul tema è emerso costantemente tra il novembre 2022 e il febbraio 2026, sviluppandosi all'interno di svariate sedute, tra cui le numero 88, 100, 234 e 314. La discussione ha interessato provvedimenti chiave come le mozioni sulle iniziative in materia energetica, il Piano Nazionale Integrato per l'Energia e il Clima (PNIEC) e i decreti-legge per la sicurezza energetica del Paese, registrando il contributo attivo di decine di deputati. Nei numerosi interventi analizzati si delinea uno scenario polarizzato sulle future strategie di approvvigionamento della Nazione.\\[6pt] Posizione del Governo\\[6pt] Per il Governo, Pichetto Fratin sostiene la necessità di sviluppare i reattori modulari per garantire competitività alle imprese e centrare gli obiettivi di decarbonizzazione. Pichetto Fratin (Forza Italia - Berlusconi Presidente - PPE, 18/06/2025) dichiara che «il nucleare contribuirà, a questo punto, a calmierare i prezzi e a ridurre la relativa volatilità».\\[6pt] Posizioni della Maggioranza\\[6pt] Per Fratelli d'Italia, Mattia sostiene l'urgenza di limitare la dipendenza tecnologica da Nazioni extra europee ripristinando le alternative atomiche. Zucconi (Fratelli d'Italia, 09/05/2023) propone di implementare la produzione energetica valutando «il ricorso al nucleare pulito e sicuro».\\[6pt] Per la Lega - Salvini Premier, Comaroli sostiene l'imprescindibilità del nucleare di ultima generazione per sopperire all'intermittenza strutturale del fotovoltaico e dell'eolico. Zinzi (Lega - Salvini Premier, 09/05/2023) ribadisce che «all'abbandono dell'energia nucleare non ha fatto seguito un piano energetico alternativo a lungo termine».\\[6pt] Per Forza Italia - Berlusconi Presidente - PPE, Squeri sostiene che l'obiettivo europeo della neutralità climatica risulti irrealizzabile senza l'apporto energetico della fissione. Squeri (Forza Italia - Berlusconi Presidente - PPE, 09/05/2023) rileva come «l'energia nucleare non è alternativa al rinnovabile, non è concorrente, anzi, semmai è concorrente al fossile».\\[6pt] Per Noi Moderati (Noi con l'Italia, Coraggio Italia, UDC e Italia al Centro) - MAIE - Centro Popolare, Cavo sostiene la necessità di investire in ricerca e nuove tecnologie per costruire una solida filiera nazionale. Colucci (Noi Moderati (Noi con l'Italia, Coraggio Italia, UDC e Italia al Centro) - MAIE - Centro Popolare, 13/12/2023) afferma che «per abbandonare la produzione di energia da fonti fossili è necessario un mix energetico, che vuol dire energia prodotta da fonti rinnovabili ed energia prodotta con il nucleare».\\[6pt] Posizioni dell'Opposizione\\[6pt] Per il Partito Democratico - Italia Democratica e Progressista, Di Sanzo sostiene che sia un grave rischio economico e strategico puntare sui nuovi reattori, ignorando le pesanti eredità del passato e i problemi di stoccaggio. Di Sanzo (Partito Democratico - Italia Democratica e Progressista, 17/04/2023) contesta le strategie della maggioranza poiché «sarebbe un grave errore portare l'Italia nuovamente nel dibattito sul nucleare di terza generazione avanzata».\\[6pt] Per il Movimento 5 Stelle, Cappelletti sostiene l'assoluta irrazionalità di una tecnologia estremamente costosa che andrebbe unicamente ad aggravare i rincari in bolletta per i cittadini. Pavanelli (Movimento 5 Stelle, 17/04/2023) denuncia il fatto che «chi oggi parla di nucleare pulito, semplicemente, non conosce i rischi che ci sono dietro il nucleare».\\[6pt] Per Alleanza Verdi e Sinistra, Zaratti sostiene l'importanza di rispettare i due referendum abrogativi passati e le irrisolte criticità legate alla gestione in sicurezza delle scorie. Bonelli (Alleanza Verdi e Sinistra, 09/05/2023) evidenzia che «il nucleare è l'esempio di un modello energetico che accumula una concentrazione di risorse pubbliche, ma le sottrae ai cittadini».\\[6pt] Per Azione - Popolari Europeisti Riformatori - Renew Europe, Benzoni sostiene che l'atomo sia l'unica energia pulita e continuativa in grado di abbassare realmente i costi di distribuzione per le imprese. Ruffino (Azione - Popolari Europeisti Riformatori - Renew Europe, 06/11/2024) auspica l'impegno istituzionale «ad aderire a qualunque dichiarazione sottoscritta tra le parti, anche in modo non unanime, che favorisca lo sviluppo globale dell'energia nucleare».\\[6pt] Per Italia Viva - Il Centro - Renew Europe, Del Barba sostiene la validità strutturale del ritorno all'atomo, criticando aspramente l'immobilismo operativo ed esecutivo dell'attuale Governo sul tema. Del Barba (Italia Viva - Il Centro - Renew Europe, 26/06/2024) osserva che «i tempi sono maturi perché si torni a studiare e a valutare, dentro il nostro PNIEC, quale potrà essere l'apporto portato dall'energia nucleare».\\[6pt] Nei documenti disponibili non emergono interventi specifici di Misto sul tema. & \\
\midrule
Salario minimo & \#\# Introduzione\\[6pt] Il dibattito parlamexwntare sull'introduzione del salario minimo si è sviluppato attraverso l'esame di diversi provvedimenti, tra cui le mozioni n. 1-00012 e abbinate e la proposta di legge A.C. 1275 recante Disposizioni per l'istituzione del salario minimo. Le discussioni hanno coinvolto decine di deputati in molteplici sedute della Camera, concentrate in particolare tra la seduta n. 15 del 29/11/2022 e la seduta n. 210 del 06/12/2023. I documenti esaminati riportano oltre quaranta interventi diretti volti a definire il perimetro delle tutele salariali nel mercato del lavoro italiano.\\[6pt] Posizione del Governo\\[6pt] Per il Governo, la Ministra del Lavoro e delle politiche sociali Calderone chiarisce che la determinazione di una retribuzione adeguata non deve avvenire tramite interventi legislativi rigidi. A tal proposito, Calderone (Misto, 11/01/2023) osserva che «Nel nostro ordinamento, la determinazione di una adeguata retribuzione non è oggi rimessa alla legge, ma è demandata alla libera negoziazione delle parti sociali attraverso lo strumento della contrattazione collettiva».\\[6pt] Posizioni della Maggioranza\\[6pt] Per Fratelli d'Italia, Malagola ritiene che l'imposizione di una soglia legale generale rischierebbe di deprimere la crescita economica e salariale complessiva, ignorando il ruolo storico dei sindacati. Nello specifico, Malagola (Fratelli d'Italia, 29/11/2022) afferma che «Il salario minimo legale rappresenta, infatti, una scorciatoia, e, come l'esperienza ci insegna, le scorciatoie portano spesso a perdere l'orientamento».\\[6pt] Per la Lega - Salvini Premier, Nisini difende l'attuale sistema italiano, sottolineando come le direttive comunitarie non obblighino all'adozione di una misura fissa ma valorizzino gli accordi di settore. Infatti, Nisini (Lega - Salvini Premier, 30/11/2022) rileva che «la contrattazione collettiva - non lo dice la Lega e non lo dice il centrodestra - nel nostro Paese tutela il 97 per cento dei lavoratori».\\[6pt] Per Forza Italia - Berlusconi Presidente - PPE, Tenerini avverte sui pericoli legati all'imposizione di una cifra valida per tutti i settori, che andrebbe inevitabilmente ad assorbire altri benefici e tutele. Nel dettaglio, Tenerini (Forza Italia - Berlusconi Presidente - PPE, 27/07/2023) contesta che «L'imposizione di un valore minimo numerico trasversale non è coerente con l'autonomia della contrattazione e rischia di essere, addirittura, sostitutiva della contrattazione stessa».\\[6pt] Per Noi Moderati (Noi con l'Italia, Coraggio Italia, UDC e Italia al Centro) - MAIE - Centro Popolare, Colucci promuove l'ampliamento dei contratti tradizionali per tutelare i lavoratori senza ricorrere a cifre imposte per legge. In Aula, Colucci (Noi Moderati (Noi con l'Italia, Coraggio Italia, UDC e Italia al Centro) - MAIE - Centro Popolare, 06/12/2023) evidenzia che «Crediamo che la strada sia la contrattazione collettiva, perché significa coinvolgimento delle parti sociali, significa sistemi di garanzia e significa combattere i contratti pirata».\\[6pt] Posizioni dell'Opposizione\\[6pt] Per il Partito Democratico - Italia Democratica e Progressista, Scotto guarda al contesto internazionale per dimostrare la compatibilità tra soglie minime di legge e solide relazioni sindacali. A supporto di ciò, Scotto (Partito Democratico - Italia Democratica e Progressista, 14/07/2025) propone che «se 22 Paesi su 27 hanno il salario minimo legale, non è perché sono impazziti, non è perché sono contrari alla contrattazione collettiva».\\[6pt] Per il Movimento 5 Stelle, Bruno si focalizza sull'urgenza di contrastare la povertà lavorativa definendo una soglia di dignità inderogabile a 9 euro. Bruno (Movimento 5 Stelle, 04/12/2023) denuncia che «Non si può accettare che passi il messaggio che per il Parlamento italiano sia giusto che in Italia si lavori per meno di 9 euro l'ora».\\[6pt] Per l'Alleanza Verdi e Sinistra, Grimaldi invoca cifre ancora più ambiziose per contrastare la perdita del potere d'acquisto dovuta all'aumento costante del costo della vita. In particolare, Grimaldi (Alleanza Verdi e Sinistra, 29/11/2022) auspica che «nessuno in Italia dovrebbe lavorare per meno di 10 euro all'ora, anzi permettetemi di dire che l'attuale situazione dell'inflazione ci dice che dovrebbero essere già 11 euro».\\[6pt] Per Azione - Popolari Europeisti Riformatori - Renew Europe, D'Alessio appoggia l'istituzione di un limite numerico sotto il quale non è possibile scendere per salvaguardare il dettato costituzionale. Pertanto, D'Alessio (Azione - Popolari Europeisti Riformatori - Renew Europe, 19/06/2024) sostiene che «C'è un salario sotto il quale - secondo noi è 9 euro - non si garantisce un'esistenza libera e dignitosa ai lavoratori».\\[6pt] Per Italia Viva - Il Centro - Renew Europe, Faraone si dice favorevole all'impianto della misura ma pone veti precisi sulle modalità di copertura finanziaria per non gravare sul ceto produttivo. Sul punto, Faraone (Italia Viva - Il Centro - Renew Europe, 01/10/2024) ribadisce che «siamo convinti del fatto il salario minimo serva in questo Paese, diciamo anche che non può essere fatto pagare con la fiscalità generale».\\[6pt] Per il Misto, Magi interpreta la fissazione del salario legale come uno strumento indispensabile per proteggere in modo mirato gli inquadramenti più vulnerabili dall'aumento dei prezzi. Precisamente, Magi (Misto, 27/07/2023) dichiara che «un salario minimo serve a far crescere i salari degli inquadramenti più bassi, più che proporzionalmente rispetto agli altri livelli salariali». & \\
\midrule
Sgravi sul costo del lavoro & \#\# Introduzione\\[6pt] Il dibattito parlamentare sul tema degli sgravi sul costo del lavoro si sviluppa attraverso l'esame di svariati disegni di legge economici e decreti in materia fiscale, coprendo un arco temporale che spazia dalla fine del 2022 ai primi mesi del 2026. L'analisi si fonda su decine di interventi in Aula che hanno coinvolto numerosi deputati afferenti a tutti gli schieramenti, nel corso di specifiche sedute come le numero 26, 126, 136, 182, 392 e 588. Le discussioni si concentrano sugli impatti macroeconomici, sull'equità delle misure e sulla sostenibilità a lungo termine delle politiche di decontribuzione.\\[6pt] Posizione del Governo\\[6pt] Per il Governo, il Ministro dell'Economia e delle finanze Giancarlo Giorgetti difende la continuità dell'Esecutivo nel ridurre progressivamente la pressione sulle buste paga. Nello specifico, Giorgetti (Lega - Salvini Premier, 22/10/2025) evidenzia che «la legge di bilancio per il 2026 proseguirà nella riduzione del prelievo fiscale sulle famiglie, estendendo i benefici finora mirati alla platea dei redditi medio-bassi ai contribuenti con redditi medi».\\[6pt] Posizioni della Maggioranza\\[6pt] Per Fratelli d'Italia, l'intervento si concentra sull'impatto positivo e misurabile degli sconti contributivi erogati direttamente a vantaggio dei dipendenti. Nel dettaglio, Schifone (Fratelli d'Italia, 10/04/2024) ribadisce che «abbiamo tagliato 6 e 7 punti portando fino a 100 euro in più nelle tasche dei lavoratori».\\[6pt] Per la Lega - Salvini Premier, si pone l'accento sull'ampiezza della platea raggiunta dalle misure di riduzione del cuneo implementate nelle manovre finanziarie. Infatti, Ottaviani (Lega - Salvini Premier, 28/12/2023) rileva che «stiamo parlando del 7 per cento per quanto riguarda i redditi fino a 25.000 euro e del 6 per cento per i redditi fino a 35.000 euro».\\[6pt] Per Forza Italia - Berlusconi Presidente - PPE, la priorità è garantire che le riduzioni del costo del lavoro acquisiscano un orizzonte stabile e duraturo per favorire la programmazione aziendale e l'occupazione. A questo proposito, Tenerini (Forza Italia - Berlusconi Presidente - PPE, 29/06/2023) auspica che «questa misura venga rinnovata nella prossima legge di bilancio e diventi strutturale».\\[6pt] Per Noi Moderati (Noi con l'Italia, Coraggio Italia, UDC e Italia al Centro) - MAIE - Centro Popolare, si sottolinea la centralità di questi sgravi in un periodo di grave difficoltà economica a causa dei picchi inflattivi. A tal proposito, Colucci (Noi Moderati (Noi con l'Italia, Coraggio Italia, UDC e Italia al Centro) - MAIE - Centro Popolare, 28/09/2023) afferma che «è già importante il taglio del cuneo fiscale che ha raggiunto per i redditi fino a 35.000 euro 5 punti e per i redditi fino a 25.000 euro 7 punti».\\[6pt] Posizioni dell'Opposizione\\[6pt] Per il Partito Democratico - Italia Democratica e Progressista, le percentuali di sgravio adottate sono considerate del tutto inadeguate a produrre miglioramenti concreti e strutturali per i dipendenti. Su questo tema, Guerra (Partito Democratico - Italia Democratica e Progressista, 22/12/2022) denuncia che «il taglio del cuneo fiscale per loro è del tutto insufficiente: pochi punti percentuali che non garantiscono alcun effetto reale».\\[6pt] Per il Movimento 5 Stelle, il taglio dei contributi deve essere utilizzato come strumento mirato e condizionato per favorire una radicale rimodulazione dei tempi di lavoro nelle aziende. Nello specifico, Aiello (Movimento 5 Stelle, 28/10/2024) propone che si applichi «un esonero contributivo - pari al 30 per cento - dei contributi previdenziali complessivi» per le imprese che adottano contratti collettivi per la riduzione dell'orario a parità di salario.\\[6pt] Nei documenti disponibili non emergono interventi specifici di Alleanza Verdi e Sinistra sul tema.\\[6pt] Per Azione - Popolari Europeisti Riformatori - Renew Europe, la necessità di incrementare concretamente le buste paga dei lavoratori più giovani passa imprescindibilmente da una forte detassazione. In quest'ottica, Bonetti (Azione - Popolari Europeisti Riformatori - Renew Europe, 22/07/2025) osserva che «si usi il taglio dell'Irpef, si usi il taglio della contribuzione, si decida tecnicamente cosa è meglio, ma vanno aumentati i salari dei giovani».\\[6pt] Per Italia Viva - Il Centro - Renew Europe, si criticano aspramente le variazioni tecniche apportate alla natura della misura che avrebbero gravemente penalizzato specifiche fasce di reddito. Nel dettaglio, Giachetti (Italia Viva - Il Centro - Renew Europe, 10/03/2025) contesta che «siete voluti intervenire snaturando, con la legge n. 207 del 2024, il valore della misura, trasformando il taglio contributivo in fiscale».\\[6pt] Per il gruppo Misto, l'appoggio alla decurtazione del costo del lavoro è strettamente condizionato all'inserimento di meccanismi normativi che favoriscano la stipula di nuovi accordi sindacali. A tal fine, Magi (Misto, 27/07/2023) precisa che «in merito al taglio del cuneo fiscale noi siamo favorevoli, a patto che questo sia di incentivo ai rinnovi contrattuali». & \\
\midrule
Golden Power & \#\# Introduzione\\[6pt] Il dibattito parlamentare sull'applicazione e l'estensione dei poteri speciali del Governo si è sviluppato trasversalmente tra il gennaio 2023 e il gennaio 2026 attraverso l'esame di decreti sull'interesse nazionale e informative sul sistema finanziario. Nei documenti analizzati si registrano gli interventi di oltre quindici deputati sul tema, distribuiti in diverse sedute plenarie tra cui la numero 44, la 520 e la 587.\\[6pt] \#\# Posizione del Governo\\[6pt] Per il Governo, il Ministro dell'Economia e delle finanze Giorgetti difende la validità dei poteri speciali per mitigare i rischi nel mercato del credito ed evitare ingerenze estere sfavorevoli. Nel dettaglio della vicenda UniCredit, Giorgetti (Lega - Salvini Premier, 30/07/2025) ribadisce che «noi non abbiamo difeso una banca, abbiamo difeso l'interesse nazionale e quindi lo strumento del golden power».\\[6pt] \#\# Posizioni della Maggioranza\\[6pt] Per Fratelli d'Italia, Rossi giudica l'attivazione dei poteri governativi uno strumento essenziale per salvaguardare il patrimonio aziendale italiano da aggressioni economiche internazionali. In merito ai limiti oggettivi di tale prerogativa, Zucconi (Fratelli d'Italia, 14/01/2026) afferma che «il golden power interverrà solo se la protezione degli interessi essenziali dello Stato non sia già garantita dalla regolamentazione di settore». Per la Lega - Salvini Premier, Andreuzza sostiene l'urgenza di rafforzare le difese istituzionali per assicurare la continuità produttiva nelle filiere maggiormente esposte alla crisi. Illustrando il perimetro di questa espansione, Andreuzza (Lega - Salvini Premier, 13/01/2026) evidenzia che le modifiche coinvolgono «la normativa in materia di poteri speciali inerenti agli attivi strategici - cosiddetto golden power - nei settori dell'energia, dei trasporti e delle comunicazioni». Per Forza Italia - Berlusconi Presidente - PPE, Bicchielli valuta positivamente i controlli preventivi per scongiurare acquisizioni ostili nelle reti strategiche del Paese senza frenare investimenti sani. In riferimento agli aiuti per le aziende coinvolte, Squeri (Forza Italia - Berlusconi Presidente - PPE, 27/01/2023) rileva che la norma introduce «la possibilità di attivare interventi di sostegno alle imprese destinatarie di misure inerenti all'esercizio dei poteri speciali riconosciuti al Governo». Per Noi Moderati (Noi con l'Italia, Coraggio Italia, UDC e Italia al Centro) - MAIE - Centro Popolare, Cavo supporta l'uso di queste difese istituzionali per impedire la delocalizzazione delle infrastrutture vitali e mantenere i livelli occupazionali dei territori. Sul tema della sovranità economica, Semenzato (Noi Moderati (Noi con l'Italia, Coraggio Italia, UDC e Italia al Centro) - MAIE - Centro Popolare, 14/01/2026) dichiara che «uno Stato che non difende gli asset strategici non è aperto, ma è vulnerabile». \\[6pt] \#\# Posizioni dell'Opposizione\\[6pt] Per il Partito Democratico - Italia Democratica e Progressista, Peluffo disapprova le logiche di mercato prive di rigorosi vincoli prescrittivi a tutela dell'ambiente e dei lavoratori in caso di passaggi di proprietà. Concentrandosi invece sul settore finanziario, Merola (Partito Democratico - Italia Democratica e Progressista, 18/12/2025) denuncia che lo strumento viene usato per il «vostro interesse di parte, per riequilibrare territorialmente a favore di altre banche e azionisti il mercato». Per il Movimento 5 Stelle, Ferrara esprime forte preoccupazione per le lacune legislative nell'economia spaziale e per le ingerenze politiche nei grandi gruppi di investimento. Circa l'impatto sulle crisi industriali, Ferrara (Movimento 5 Stelle, 15/11/2024) propone al Governo di «farsi garante dell'applicazione rigorosa del golden power non per autorizzare ma per condizionare e bloccare dove serve». Nei documenti disponibili non emergono interventi specifici di Alleanza Verdi e Sinistra sul tema. Per Azione - Popolari Europeisti Riformatori - Renew Europe, Richetti biasima l'assenza di stringenti clausole di salvaguardia per la continuità lavorativa nelle recenti direttive emanate dal Consiglio dei ministri. In ambito bancario, Benzoni (Azione - Popolari Europeisti Riformatori - Renew Europe, 18/12/2025) contesta le dinamiche delle recenti acquisizioni sostenendo che «tutto è influenzato dall'applicazione della golden power iniziale del Governo». Per Italia Viva - Il Centro - Renew Europe, Del Barba condanna l'introduzione di criteri discrezionali che limitano la libera concorrenza tra attori finanziari paritari. Sul blocco delle acquisizioni nel credito, Del Barba (Italia Viva - Il Centro - Renew Europe, 18/12/2025) osserva che adoperare il meccanismo di veto tra banche nazionali vigilate rappresenta «un utilizzo quantomeno spropositato e sicuramente inedito». Per il Misto, Della Vedova attacca l'utilizzo governativo giudicato arbitrario dei poteri di sicurezza nazionale in dinamiche di legittima competizione azionaria europea. Riferendosi all'offerta pubblica nel comparto bancario, Della Vedova (Misto, 12/03/2025) auspica «che il golden power non venga usato in modo strumentale per interferire nel riequilibrio». & \\

\end{longtable}
\end{landscape}

\end{document}